
\chapter{Research Roadmap and Deployment Path}
\label{ch:roadmap}

This chapter lays out the sequence of work required to move ORIUS from a
battery-validated prototype to a multi-domain, deployable safety runtime.
Each milestone is ordered by dependency: later items require earlier ones,
and no step can be skipped without accumulating technical debt.

\section{Submit the Battery Foundation First}

The immediate priority is to publish the battery-domain results as a
self-contained contribution.  The submission package comprises the CPSBench
harness, the DC3S implementation, the dual-trajectory evaluation protocol,
and the empirical evidence of zero true-state violations.  Peer review of
the battery result establishes the OASG concept and the DC3S shield in the
literature before generalisation claims are made.

\section{Fix Remaining Empirical Gaps}

Before extending to new domains, the battery evidence must be hardened:

\begin{enumerate}
  \item \textbf{Out-of-distribution faults.}  Extend the fault model to
    include correlated multi-sensor dropout and Byzantine sensor failures.
    Validate that CQR residuals remain exchangeable under these conditions
    or document the breakdown.
  \item \textbf{Long-horizon aging.}  Obtain or simulate multi-year
    degradation trajectories and verify that the aging-aware recalibration
    of Chapter~\ref{ch:temporal-extensions} maintains coverage.
  \item \textbf{Regional breadth.}  Run the full battery on at least two
    additional grid regions (e.g., ERCOT, NordPool) to confirm that
    market-structure differences do not invalidate the safety guarantee.
\end{enumerate}

\section{Extract ORIUS Core Modules}

Refactor the current codebase to separate domain-independent modules
(OQE scoring, RAC inflation, certificate issuance, audit trail) from
battery-specific code (SOC dynamics, battery constraints, SCADA parser).
The result is a \texttt{orius-core} library with a stable API and a
\texttt{orius-battery} adapter that depends on it.  The separation must
be enforced by interface boundaries, not merely by convention.

\section{Build the Battery Adapter Package}

Package the battery-specific adapter as a standalone installable module
that depends on \texttt{orius-core}.  This package includes:

\begin{itemize}
  \item The SOC integrator plant model.
  \item The SCADA telemetry parser and OQE feature extractor.
  \item The box-constraint projection oracle.
  \item The CPSBench battery plant for regression testing.
\end{itemize}

Publishing this package demonstrates the adapter pattern in practice and
provides a template for subsequent domain adapters.

\section{Add a Second Domain}

The second domain should stress the aspects of ORIUS that battery did not
test: nonlinear dynamics, multi-dimensional state, and non-box constraints.
Candidates include building HVAC (nonlinear thermal dynamics, polytopic
comfort constraints) and autonomous vehicle lane-keeping (six-DOF state,
curvature-dependent road constraints).  The second-domain study must
reproduce the full evaluation protocol: dual-trajectory simulation, TSVR
measurement, conformal coverage audit, and intervention-rate analysis.

\section{Add a Third Domain and Beyond}

Each additional domain further validates the invariant core and reveals
edge cases in the adapter interface.  After the third domain the
\texttt{orius-core} API should be considered stable; subsequent domains
become engineering exercises rather than research contributions.

\section{Hardware-in-the-Loop Validation}

The CPSBench harness simulates telemetry faults in software.  Before
deployment, the DC3S shield must be tested with \emph{real} packet loss,
delay jitter, and sensor degradation on a hardware-in-the-loop (HIL)
testbed.  The HIL stage validates:

\begin{enumerate}
  \item That the OQE feature extractor behaves correctly on raw network
    packets, not simulated dropout masks.
  \item That the shield's latency budget (established in
    Chapter~\ref{ch:latency}) holds on embedded hardware.
  \item That the certificate chain persists correctly under power
    interruptions and storage faults.
\end{enumerate}

\section{ORIUS-Bench Public Release}

The CPSBench harness should be released as a public benchmark---ORIUS-Bench
---with a standardised evaluation API, a leaderboard, and reference
implementations for at least two domains.  Public benchmarking invites
independent reproduction and adversarial stress-testing by the community,
which no internal audit can substitute.

\section{CertOS Runtime Path}

The long-term deployment target is a \emph{CertOS}: a lightweight runtime
operating system that wraps any black-box controller with the ORIUS
safety layer.  CertOS would run as a sidecar process, intercept actuation
commands, evaluate them against the current certificate, and either pass
or project.  The runtime must satisfy hard real-time constraints, which
imposes implementation requirements beyond the Python prototype used in
this dissertation.

\section{Chapter Summary}

The roadmap proceeds in strict dependency order: publish the battery
foundation, harden its empirical gaps, extract the domain-independent
core, validate on a second and third domain, run hardware-in-the-loop
tests, release ORIUS-Bench publicly, and ultimately build a deployable
CertOS runtime.  Each step is necessary; none is sufficient alone.  The
battery result is the foundation, not the destination.
